\documentclass[10pt,letterpaper]{article}
\usepackage{cornell-notes}

\title{Chapter 7: The Application Layer}
\author{}
\date{}

\setDocTitle{Chapter 7: The Application Layer}
\setDocSubtitle{Computer Networks}
\setDocVersion{v1.0}
\setDocStatus{Study Companion}
\setDocOwner{Computer Networks Cornell Notes}
\setDocAuthor{}
\setDocDate{\today}
\setDocSummary{Cornell-format study and review notes for Chapter 7: The Application Layer.}

\setCornellCollection{Computer Networks Cornell Notes}
\setCornellUnitType{Chapter}
\setCornellUnitNumber{7}
\setCornellUnitTitle{The Application Layer}
\setCornellCompanionLabel{Study and review companion}

\hypersetup{pdftitle={Chapter 7: The Application Layer},pdfsubject={Cornell notes},pdfauthor={}}

\begin{document}
\maketitle
\notesection{Chapter overview}
\begin{overviewbox}
\textbf{Learning objectives}\begin{itemize}
  \item Trace DNS resolution through names, resource records, zones, and servers.
  \item Explain electronic-mail architecture, message formats, transfer, and final delivery.
  \item Relate URLs, HTTP, static content, dynamic applications, and Web architecture.
  \item Connect audio/video representation to stored streaming, live streaming, and conferencing constraints.
  \item Compare server farms, proxies, content-delivery networks, and peer-to-peer distribution.
\end{itemize}
\textbf{Prerequisite concepts}\quad TCP and UDP services, sockets, hierarchical naming, digital sampling, caching, and basic security concepts.\par
\textbf{Key themes}\quad Naming; request/response; store-and-forward messaging; representation; caching; timing; replication; edge distribution.\par
\textbf{Named mechanisms and standards}\quad DNS, SMTP, MIME, POP3, IMAP, URLs, HTTP, HTML, cookies, server-side and client-side processing, audio/video coding, RTP-related media transport, CDNs, and peer-to-peer systems.
\end{overviewbox}

\notesection{Cornell cue-and-note record}

\begin{longtable}{@{}L{0.245\textwidth}|L{0.695\textwidth}@{}}
\rowcolor{CNavy}\color{white}\bfseries Cue / recall prompt &
\color{white}\bfseries Notes \\
\endfirsthead
\rowcolor{CNavy}\color{white}\bfseries Cue / recall prompt &
\color{white}\bfseries Notes (continued) \\
\endhead
\hline
\endfoot
\cellcolor{CCue}\textbf{7.1 DNS} & \begin{minipage}[t]{\linewidth}\raggedright DNS maps structured names to resource data through a distributed, delegated, cached database.\par\begin{itemize}\item \textbf{7.1.1 DNS name space} The hierarchical name space is divided into domains and zones. Labels form domain names; administrative delegation lets different authorities manage subtrees without one central database.\item \textbf{7.1.2 Resource records} Records associate names with typed data such as address, name-server, mail-exchanger, alias, and authority information. A time-to-live bounds how long a cached answer should be retained.\item \textbf{7.1.3 Name servers} Resolvers query recursive or iterative servers. Root, top-level-domain, and authoritative servers guide resolution; caching reduces latency and load but permits temporarily stale results.\end{itemize}\end{minipage} \\ \hline
\cellcolor{CCue}\textbf{7.2 Electronic mail} & \begin{minipage}[t]{\linewidth}\raggedright Electronic mail is an asynchronous store-and-forward application, not a single end-to-end protocol.\par\begin{itemize}\item \textbf{7.2.1 Architecture and services} User agents submit messages; transfer agents relay them across mail infrastructure; delivery agents place them in recipient mailboxes. Mailing lists, aliases, and gateways add distribution and translation.\item \textbf{7.2.2 User agent} The user agent composes, addresses, displays, replies to, forwards, files, and searches messages while interacting with submission and mailbox services.\item \textbf{7.2.3 Message formats} Internet messages contain an envelope used for transport and a header/body representation for users. MIME describes content types, transfer encodings, and multipart bodies beyond plain ASCII text.\item \textbf{7.2.4 Message transfer} SMTP moves mail through command/response exchanges over reliable transport. Relays use destination information and queue messages when immediate delivery is impossible.\item \textbf{7.2.5 Final delivery} A local delivery agent writes the mailbox. POP3 primarily supports download-oriented access; IMAP supports server-resident folders, state, and selective synchronization. Webmail presents mailbox access through a Web application.\end{itemize}\end{minipage} \\ \hline
\cellcolor{CCue}\textbf{7.3 World Wide Web} & \begin{minipage}[t]{\linewidth}\raggedright \begin{itemize}\item \textbf{7.3.1 Architectural overview} Browsers identify resources with URLs, request representations from servers, interpret content, and follow embedded links. Caches and proxies can serve or forward requests between clients and origin servers.\item \textbf{7.3.2 Static Web pages} HTML describes document structure and links; style and embedded objects affect presentation. A static resource is returned from stored content without generating a response from request-specific server computation.\item \textbf{7.3.3 Dynamic pages and applications} Server-side programs generate representations from request data and backend state; client-side code changes the user interface and exchanges data after the initial page. State may be carried by URLs, hidden fields, cookies, or server sessions.\item \textbf{7.3.4 HTTP} HTTP uses request methods, resource identifiers, headers, status codes, and message bodies. Persistent connections reduce setup overhead; caching and conditional requests avoid unnecessary transfers. The protocol is stateless at the request layer even when applications maintain sessions.\item \textbf{7.3.5 Mobile Web} Small screens, variable radio capacity, latency, device capability, and energy motivate adaptation, compact representations, and careful application design.\item \textbf{7.3.6 Web search} Search engines crawl links, index content, rank candidate pages, and answer queries. Crawling coverage, duplicate content, ranking signals, and adversarial manipulation affect results.\end{itemize}\end{minipage} \\ \hline
\cellcolor{CCue}\textbf{7.4 Streaming audio and video} & \begin{minipage}[t]{\linewidth}\raggedright \begin{itemize}\item \textbf{7.4.1 Digital audio} Sampling converts amplitude measurements into numbers. Sampling rate, quantization precision, channels, and compression determine quality and bit rate; lossy coding removes information judged less perceptible.\item \textbf{7.4.2 Digital video} Video is a sequence of images whose spatial and temporal redundancy can be compressed. Intra-frame coding compresses one image; inter-frame prediction describes change between images.\item \textbf{7.4.3 Stored-media streaming} The receiver begins playback before the complete object arrives. Buffering absorbs network jitter; startup delay and playout rate must be balanced against stalls.\item \textbf{7.4.4 Live-media streaming} Content is produced as it is sent, so the system cannot prefetch arbitrarily far ahead. Buffering adds latency while smoothing jitter; multicast can reduce duplicate network traffic.\item \textbf{7.4.5 Real-time conferencing} Interactive audio/video has strict delay needs, so late packets may be less useful than lost packets. Adaptive coding, timestamps, sequence information, jitter buffers, and feedback support conversation quality.\end{itemize}\end{minipage} \\ \hline
\cellcolor{CCue}\textbf{7.5 Content delivery} & \begin{minipage}[t]{\linewidth}\raggedright \begin{itemize}\item \textbf{7.5.1 Content and Internet traffic} Popular and large objects create repeated transfers and hotspots. Locality and skewed popularity make caching and replication effective.\item \textbf{7.5.2 Server farms and Web proxies} A front end distributes requests among replicated servers; proxies cache on behalf of clients or organizations. Load balancing, cache consistency, and failure handling affect results.\item \textbf{7.5.3 Content-delivery networks} A CDN places replicas near users and directs requests to a suitable edge server using naming or application redirection. The chosen server reflects location, load, health, and content availability.\item \textbf{7.5.4 Peer-to-peer networks} Peers contribute upload, storage, and discovery capacity. Centralized indexes, flooding, distributed hash tables, and swarming offer different location and distribution tradeoffs; incentives affect cooperation.\end{itemize}\end{minipage} \\ \hline
\cellcolor{CCue}\textbf{7.6 Summary} & \begin{minipage}[t]{\linewidth}\raggedright Application protocols transform transport services into naming, mail, Web, and media functions. Caching, replication, adaptation, and distribution improve scale and performance but introduce staleness, consistency, placement, and trust concerns.\end{minipage} \\ \hline
\end{longtable}

\begin{legacybox}The presented Web, mobile-Web, mail, codecs, search, and peer-to-peer examples reflect the period of the source. Protocol principles are preserved, but specific browser techniques and deployment patterns should be treated as historical or legacy context where indicated.\end{legacybox}

\notesection{Terminology and notation}

\begin{longtable}{@{}L{0.245\textwidth}|L{0.695\textwidth}@{}}
\rowcolor{CNavy}\color{white}\bfseries Cue / recall prompt &
\color{white}\bfseries Notes \\
\endfirsthead
\rowcolor{CNavy}\color{white}\bfseries Cue / recall prompt &
\color{white}\bfseries Notes (continued) \\
\endhead
\hline
\endfoot
\cellcolor{CCue}\textbf{Domain name} & \begin{minipage}[t]{\linewidth}\raggedright Hierarchical sequence of DNS labels naming a point in the global name tree.\end{minipage} \\ \hline
\cellcolor{CCue}\textbf{Zone} & \begin{minipage}[t]{\linewidth}\raggedright Administratively managed portion of the DNS name space served by authoritative data.\end{minipage} \\ \hline
\cellcolor{CCue}\textbf{Resource record} & \begin{minipage}[t]{\linewidth}\raggedright Typed DNS data associated with a name and a cache lifetime.\end{minipage} \\ \hline
\cellcolor{CCue}\textbf{MIME} & \begin{minipage}[t]{\linewidth}\raggedright Message format extensions for typed, encoded, and multipart electronic-mail content.\end{minipage} \\ \hline
\cellcolor{CCue}\textbf{URL} & \begin{minipage}[t]{\linewidth}\raggedright Identifier that specifies how and where a Web resource is addressed.\end{minipage} \\ \hline
\cellcolor{CCue}\textbf{HTTP} & \begin{minipage}[t]{\linewidth}\raggedright Application protocol using requests and responses to transfer resource representations and control information.\end{minipage} \\ \hline
\cellcolor{CCue}\textbf{Cookie} & \begin{minipage}[t]{\linewidth}\raggedright Small name/value state carried between a user agent and server under defined scope rules.\end{minipage} \\ \hline
\cellcolor{CCue}\textbf{Jitter buffer} & \begin{minipage}[t]{\linewidth}\raggedright Receiver buffer that trades additional playout delay for tolerance of variable packet arrival time.\end{minipage} \\ \hline
\cellcolor{CCue}\textbf{CDN} & \begin{minipage}[t]{\linewidth}\raggedright Distributed content-delivery system that serves replicated objects from selected edge locations.\end{minipage} \\ \hline
\cellcolor{CCue}\textbf{DHT} & \begin{minipage}[t]{\linewidth}\raggedright Distributed hash table that maps keys to responsible peers without one central index.\end{minipage} \\ \hline
\end{longtable}

\notesection{Comparison matrix}

\begin{longtable}{@{}L{0.313\textwidth}L{0.313\textwidth}L{0.313\textwidth}@{}}
\toprule
\textbf{Mechanism} & \textbf{Primary goal} & \textbf{Key limitation / cost} \\ \midrule
DNS caching & Reduce lookup delay and authoritative load & Answers may remain stale until lifetime expires \\
POP3 & Simple mailbox download & Limited server-side organization and synchronization \\
IMAP & Manage server-resident mailbox state & More server state and protocol complexity \\
Forward Web proxy & Serve/cache requests for clients & Visibility, policy, and cache correctness \\
CDN edge & Replicate content near users & Placement, redirection, consistency, and operational cost \\
Peer-to-peer swarm & Use peer upload capacity & Discovery, churn, incentives, and trust \\
\bottomrule
\end{longtable}
\notesection{Chapter synthesis}

\begin{overviewbox}\textbf{Cause-and-effect chain.} Human-friendly names require distributed resolution; asynchronous messaging requires queues and relays; linked resources require identifiers and request/response transfer; media timing requires coding, buffering, and playout; global demand requires caches and replicas near consumers.\par
\textbf{Common confusions.} DNS is not a general directory of all application state; SMTP transfers mail while POP3/IMAP access mailboxes; HTTP statelessness does not prevent application sessions; downloading is not the same as streaming; and a CDN differs from a client's forward proxy.\par
\textbf{Derived insight.} Many application-layer scaling techniques deliberately serve a copy rather than the origin, turning placement and freshness into first-class protocol concerns.\end{overviewbox}

\notesection{Retrieval practice}

\begin{enumerate}
  \item Trace iterative DNS resolution for an uncached multi-label name.
  \item Why do DNS records have time-to-live values?
  \item Separate the roles of a mail user agent, transfer agent, and delivery agent.
  \item What problems do MIME content types and transfer encodings solve?
  \item Compare SMTP, POP3, and IMAP by direction and responsibility.
  \item Trace an HTTP request through a cache to an origin server and back.
  \item How can an application maintain a session over stateless HTTP requests?
  \item Distinguish static content, server-generated content, and client-side application behavior.
  \item Relate sampling rate, quantization, and compression to digital-audio bit rate and quality.
  \item Why does a jitter buffer improve continuity while increasing latency?
  \item Compare stored streaming, live streaming, and interactive conferencing constraints.
  \item How does a CDN choose an edge, and why might the geographically closest edge not be best?
  \item Compare centralized peer indexes, flooding, and distributed hash tables.
  \item Why do incentive mechanisms matter in peer-to-peer distribution?
\end{enumerate}
\begin{CornellSummaryBox}{Rapid-review Cornell summary}DNS resolves delegated names through typed records and caches. Mail separates composition, transfer, and mailbox access. The Web combines URLs, HTTP, representations, and application state. Streaming adds timing and buffering requirements. Server farms, proxies, CDNs, and peers replicate work and content to scale delivery.\end{CornellSummaryBox}
\end{document}
